% c11.fig.bereiche — the three sets of functions in the Fundament block, as nested regions.
%
% A schematic, not a plot.  It restates the inclusions used with Theorem X.15 of Reed-Simon II
% (p. 165, read from the page image) and Corollary c11.cor.relativ:
%     C_c^infty(R^3)  subset  D(-Lap)  subset  L^2(R^3);
% H_rel is essentially self-adjoint on the innermost set and self-adjoint on the middle one.
% Greyscale: outlines and labels only.
\begin{tikzpicture}[x=1cm, y=1cm]
  \draw[tpc axis, rounded corners=6pt] (0,0) rectangle (9.6,3.6);
  \draw[tpc axis, rounded corners=6pt] (0.5,0.35) rectangle (7.4,2.75);
  \draw[tpc axis, rounded corners=6pt, pattern=north east lines, pattern color=tpc-plate-rule] (1.0,0.7) rectangle (4.3,1.9);
  \node[tpc label, anchor=north east] at (9.4,3.45) {$L^2(\mathbb{R}^3)$};
  \node[tpc label, anchor=north east] at (7.2,2.6) {$D(-\Delta)$};
  \node[tpc label, fill=tpc-plate, inner sep=2pt] at (2.65,1.3) {$C_c^{\infty}(\mathbb{R}^3)$};
  \node[tpc label, anchor=west, align=left] at (4.5,1.05)
    {\tpcifenglish{$\op{H_{\mathrm{rel}}}$ essentially}{$\op{H_{\mathrm{rel}}}$ wesentlich}\\
     \tpcifenglish{self-adjoint here}{selbstadjungiert hier}};
  \draw[tpc leader] (4.5,1.05) -- (4.3,1.2);
  \node[tpc label, anchor=south west, align=left] at (0.8,2.05)
    {\tpcifenglish{$\op{H_{\mathrm{rel}}}$ self-adjoint here}{$\op{H_{\mathrm{rel}}}$ selbstadjungiert hier}};
\end{tikzpicture}
